% from ega3-1-fr.tex
\oldpage[III]{105}
\[
s^*(\mathcal O_{X''}(n))=\mathcal O_{X'}(nd).
\]
On établira donc l'existence d'un $Y$-isomorphisme $X\xrightarrow{\sim}X'$
si l'on prouve la

\begin{env}[2.3.4.1]
\label{III.2.3.4.1-fr}
Soient $Y$ un préschéma noethérien, $q:X\to Y$ un morphisme projectif,
$\mathcal L$ un $\mathcal O_X$-Module inversible très ample pour $q$.
Alors $\mathcal S=\bigoplus_{n\geq0}q_*(\mathcal L^{\otimes n})$ est une
$\mathcal O_Y$-Algèbre graduée quasi-cohérente de type fini, telle que
$\mathcal S_n=\mathcal S_1^n$ pour $n$ assez grand, et il existe un
$Y$-isomorphisme $r:X\xrightarrow{\sim}P=\operatorname{Proj}(\mathcal S)$
tel que $\mathcal L=r^*(\mathcal O_P(1))$.
\end{env}

Comme $q$ est un morphisme projectif, il résulte de (II, 5.4.4 et 4.4.7)
qu'il existe un $Y$-isomorphisme
$r':X\xrightarrow{\sim}P'=\operatorname{Proj}(\mathcal T)$, où
$\mathcal T$ est une $\mathcal O_Y$-Algèbre quasi-cohérente,
$\mathcal T_1$ est de type fini et engendre $\mathcal T$, et
$\mathcal L=r'^*(\mathcal O_{P'}(1))$. Si $q':P'\to Y$ est le morphisme
structural, on a
\[
\mathcal S=\bigoplus_{n\geq0}q'_*(\mathcal O_{P'}(n)).
\]
En vertu de (\hyperref[III.2.3.1-fr]{2.3.1}), l'homomorphisme canonique
$\alpha_n:\mathcal T_n\to\mathcal S_n=q'_*(\mathcal O_{P'}(n))$ est
bijectif pour $n$ assez grand~; comme $\mathcal T_n=\mathcal T_1^n$, on a
\emph{a fortiori} $\mathcal S_n=\mathcal S_1^n$ dès que $n$ est assez
grand. En outre, l'homomorphisme canonique
$\alpha:\mathcal T\to\mathcal S$ de $\mathcal O_Y$-Algèbres graduées est
un \textbf{(TN)}-isomorphisme~; par suite
\phantomsection\label{III.2.3.5.1-fr}
\[
\Phi=\operatorname{Proj}(\alpha):
\operatorname{Proj}(\mathcal S)\longrightarrow
\operatorname{Proj}(\mathcal T)
\]
est un isomorphisme (II, 3.6.1). En outre, il résulte de (II, 3.5.2) et
du fait que les $\mathcal T$-Modules gradués obtenus à partir des
$\mathcal S(n)$ par restriction des scalaires sont \textbf{(TN)}-isomorphes
aux $\mathcal T(n)$ que
\[
\Phi_*(\mathcal O_P(n))=\mathcal O_{P'}(n)
\]
pour tout $n$ (II, 3.4.2). Pour achever de prouver
(\hyperref[III.2.3.4.1-fr]{2.3.4.1}), il reste à montrer que $\mathcal S$
est une $\mathcal O_Y$-Algèbre de type fini~; or les
$\mathcal S_n=q'_*(\mathcal O_{P'}(n))$ sont des
$\mathcal O_Y$-Modules cohérents en vertu de
(\hyperref[III.2.2.1-fr]{2.2.1}), et si
$\mathcal S_n=\mathcal S_1^n$ pour $n\geq n_0$, $\mathcal S$ est
engendrée par $\bigoplus_{i\leq n_0}\mathcal S_i$, qui est cohérent~;
d'où notre assertion (I, 9.6.2).

Revenons à la démonstration de
(\hyperref[III.2.3.4-fr]{2.3.4}), dont nous reprenons les notations. Nous
avons démontré l'existence d'un $Y$-isomorphisme $r'':X\xrightarrow{\sim}X''$
tel que $r''_*(\mathcal L^{\otimes n})=\mathcal O_{X''}(n)$ pour tout
$n\in\mathbb Z$~; nous désignerons par $q''$ le morphisme structural
$X''\to Y$. Notons maintenant que $\mathcal S'$ est somme directe des
$\mathcal S'^{(d)}$-Modules gradués $\mathcal S'^{(d,k)}$~; chacun de ces
derniers est un $\mathcal S'^{(d)}$-Module quasi-cohérent de type fini, en
vertu de la périodicité de $\mathcal S'$ et de l'hypothèse que les
$\mathcal S'_n$ sont des $\mathcal O_X$-Modules de type fini
(II, 8.14.12). Posons
$\mathcal F^{(k)}=r''_*(\mathcal S'^{(d,k)})$, de sorte que les
$\mathcal F^{(k)}$ sont des $\mathcal S_{X''}$-Modules gradués
quasi-cohérents de type fini~; par suite (II, 8.14.8), l'homomorphisme
canonique
\[
\beta:\shProj(\bbGamma_*(\mathcal F^{(k)}))
\longrightarrow\mathcal F^{(k)}
\]
est un isomorphisme de $\mathcal S_{X''}$-Modules. Mais on a
$q''_*((\mathcal F^{(k)})_n)=q_*((\mathcal S'^{(d,k)})_n)$, et pour
$n\geq0$ ce dernier $\mathcal O_Y$-Module est par définition égal à
$(\mathcal S^{(d,k)})_n$. Autrement dit, l'injection canonique
$\mathcal S^{(d,k)}\to\bbGamma_*(\mathcal F^{(k)})$ est un
\textbf{(TN)}-isomorphisme~; donc
\[
\shProj(\mathcal S^{(d,k)})
=\shProj(\bbGamma_*(\mathcal F^{(k)})),
\]
et par suite
$r''{}^*(\shProj(\mathcal S^{(d,k)}))\simeq\mathcal S'^{(d,k)}$.
Il reste à remarquer que, à un isomorphisme canonique près,
\[
\shProj(\mathcal S^{(d,k)})
=s_*((\shProj(\mathcal S))^{(d,k)})
\]
(II, 8.14.13.1) pour avoir démontré l'isomorphisme de
$r^*(\shProj(\mathcal S))$ et de $\mathcal S'$.

Enfin, en vertu de (\hyperref[III.2.3.2-fr]{2.3.2}), chacun des
$\bbGamma_*(\mathcal F^{(k)})$ vérifie la condition \textbf{(TF)}, donc il
en est de même de chacun des $\mathcal S^{(d,k)}$~; en outre, les
$\mathcal S_n=q_*(\mathcal S'_n)$ sont cohérents en vertu de
(\hyperref[III.2.2.1-fr]{2.2.1}), puisque les $\mathcal S'_n$ sont
cohérents. On en conclut aussitôt que les $\mathcal S^{(d,k)}$ sont des
$\mathcal S^{(d)}$-Modules de type fini~; comme
(\hyperref[III.2.3.4.1-fr]{2.3.4.1}) montre que $\mathcal S^{(d)}$ est une
$\mathcal O_Y$-Algèbre de type fini, il en est de même de $\mathcal S$.
\end{proof}

